\section{Conclusion}

This paper audits and predicts a constructed SME fragility classification using Rwanda 2023 Enterprise Survey data. Its main contribution is not simply a high-performing predictive model. Rather, the paper shows how a legacy fragility indicator can be reconstructed, checked against official metadata, tested for predictor leakage, and evaluated with transparent validation diagnostics in a low-data African enterprise context. This index-auditing workflow is relevant beyond Rwanda because many social-science and development-policy datasets rely on constructed screening indicators whose validity depends on transparent coding, documentation, and external validation.

The substantive conclusion is cautious. The full model replicates the recovered screening rule very well, but the leakage-reduced benchmark provides the more meaningful estimate of independent predictive signal. That benchmark performs moderately, which supports continued development of the screening framework but not immediate policy deployment as a definitive risk score. The construct-validity audit further shows that the current outcome should be interpreted as an innovation, website-status, and employment-linked screening classification, not as a direct measure of observed firm failure or financial distress.

Before journal submission and policy use, the index should be validated against independent outcomes, tested with survey-design adjustments, and governed as a supportive diagnostic tool. Used carefully, the framework can help researchers and policymakers identify where more evidence is needed and how SME-support systems might better detect vulnerability before observable failure occurs.
